\chapter{Sign conventions and shifts}
\label{app:signs}

\index{sign conventions|textbf}

The sign conventions used throughout this monograph are chosen for compatibility with
Loday--Vallette~\cite{LV12} (operadic structures),
Beilinson--Drinfeld~\cite{BD04} (chiral algebras), and Lurie~\cite{HA}
($\infty$-categorical foundations).

\section{The Koszul sign rule}
\label{sec:koszul-sign-rule}

\begin{convention}[Koszul sign rule]
\label{conv:koszul}
\index{Koszul sign rule|textbf}
Let $\cC$ be a graded category (chain complexes, differential graded algebras, 
graded vector spaces, etc.). When two homogeneous elements $a$ and $b$ of 
degrees $|a|$ and $|b|$ respectively are interchanged in any formula, a sign 
$(-1)^{|a| \cdot |b|}$ must be introduced:
\[
a \otimes b \longmapsto (-1)^{|a| \cdot |b|} \, b \otimes a.
\]
This rule applies universally: to tensor products, compositions, evaluations, 
and any operation where the order of graded objects matters.
\end{convention}

\begin{definition}[Graded commutator]
\label{def:graded-commutator}
For homogeneous elements $a, b$ in a graded algebra $A$, the 
\emph{graded commutator} is:
\[
[a, b] := a \cdot b - (-1)^{|a| \cdot |b|} b \cdot a.
\]
An algebra is \emph{graded commutative} if $[a, b] = 0$ for all homogeneous 
$a, b$, which is equivalent to $a \cdot b = (-1)^{|a| \cdot |b|} b \cdot a$.
\end{definition}

\begin{proposition}[Graded Jacobi identity; \ClaimStatusProvedHere]
\label{prop:graded-jacobi}
For any graded Lie algebra $(L, [-,-])$, the bracket satisfies the 
\emph{graded Jacobi identity}:
\[
(-1)^{|a| \cdot |c|} [a, [b, c]] + (-1)^{|b| \cdot |a|} [b, [c, a]] 
+ (-1)^{|c| \cdot |b|} [c, [a, b]] = 0.
\]
This is equivalent to requiring that the adjoint action $\ad_a = [a, -]$ 
is a derivation of the bracket:
\[
[a, [b, c]] = [[a, b], c] + (-1)^{|a| \cdot |b|} [b, [a, c]].
\]
\end{proposition}

\begin{proof}
The equivalence follows from expanding the graded Jacobi identity and 
applying graded antisymmetry $[b, c] = -(-1)^{|b| \cdot |c|} [c, b]$. 
The cyclic sum becomes:
\begin{align*}
0 &= (-1)^{|a| \cdot |c|} [a, [b, c]] + (-1)^{|b| \cdot |a|} [b, [c, a]] 
    + (-1)^{|c| \cdot |b|} [c, [a, b]] \\
  &= (-1)^{|a| \cdot |c|} [a, [b, c]] 
    - (-1)^{|b| \cdot |a|} (-1)^{|a| \cdot |c|} [b, [a, c]] 
    + (-1)^{|c| \cdot |b|} [c, [a, b]] \\
  &= (-1)^{|a| \cdot |c|} \bigl( [a, [b, c]] - (-1)^{|a| \cdot |b|} [b, [a, c]] 
    - [[a, b], c] \bigr),
\end{align*}
using $[[a,b],c] = -(-1)^{(|a|+|b|)|c|}[c,[a,b]]$ and simplifying exponents.
\end{proof}

\begin{convention}[Sign in differential graded structures]
\label{conv:dg-signs}
In a differential graded algebra $(A, d, \cdot)$:
\begin{enumerate}[label=(\roman*)]
\item The differential has degree $|d| = +1$ (cohomological convention) 
or $|d| = -1$ (homological convention).
\item The Leibniz rule incorporates the Koszul sign:
\[
d(a \cdot b) = (da) \cdot b + (-1)^{|a|} a \cdot (db).
\]
\item For a differential graded Lie algebra, the bracket-differential 
compatibility is:
\[
d[a, b] = [da, b] + (-1)^{|a|} [a, db].
\]
\end{enumerate}
Throughout this monograph, we use the \emph{cohomological convention} 
$|d| = +1$ unless explicitly stated otherwise.
\end{convention}

\begin{lemma}[Sign rule for compositions \cite{LV12}; \ClaimStatusProvedElsewhere]
\label{lem:composition-signs}
Let $f: V \to W$ and $g: W \to X$ be graded linear maps of degrees $|f|$ 
and $|g|$ respectively. The composition $g \circ f: V \to X$ has degree 
$|g \circ f| = |f| + |g|$. For evaluation on $v \in V$:
\[
(g \circ f)(v) = g(f(v)) \quad \text{(no sign)}.
\]
However, when expressing $g \circ f$ in terms of tensor factors, the 
Koszul rule applies:
\[
(f \otimes g) \circ \tau_{V,W} = (-1)^{|f| \cdot |g|} \tau_{V,X} \circ (g \otimes f)
\]
where $\tau$ denotes the graded twist map.
\end{lemma}

\begin{definition}[Graded Hom complex]
\label{def:graded-hom}
For chain complexes $(V, d_V)$ and $(W, d_W)$, the \emph{internal Hom complex} 
$\underline{\Hom}(V, W)$ has:
\begin{enumerate}[label=(\roman*)]
\item Degree $n$ component: $\underline{\Hom}(V, W)^n = \prod_{k \in \bZ} 
\Hom_k(V^k, W^{k+n})$, i.e., degree-$n$ maps.
\item Differential: For $f \in \underline{\Hom}(V, W)^n$,
\[
(d_{\underline{\Hom}} f)(v) := d_W(f(v)) - (-1)^{|f|} f(d_V(v)).
\]
\end{enumerate}
The sign $(-1)^{|f|}$ ensures $d_{\underline{\Hom}}^2 = 0$.
\end{definition}

\begin{verification}[\texorpdfstring{$d_{\underline{\Hom}}^2 = 0$}{d\_Hom\textasciicircum 2 = 0}]
\label{ver:d-hom-squared}
For $f$ of degree $|f|$:
\begin{align*}
(d^2 f)(v) &= d\bigl( d_W(f(v)) - (-1)^{|f|} f(d_V v) \bigr) 
             - (-1)^{|f|+1} \bigl( d_W(f(v)) - (-1)^{|f|} f(d_V v) \bigr)|_{v \mapsto d_V v} \\
&= d_W^2(f(v)) - (-1)^{|f|} d_W(f(d_V v)) 
   - (-1)^{|f|+1} d_W(f(d_V v)) + (-1)^{2|f|+1} f(d_V^2 v) \\
&= 0 - (-1)^{|f|} d_W(f(d_V v)) + (-1)^{|f|} d_W(f(d_V v)) + 0 = 0.
\end{align*}
The cancellation depends on the sign convention.
\end{verification}


\section{Cohomological versus homological grading}
\label{sec:cohom-vs-homol}


\begin{definition}[Grading conventions]
\label{def:grading-conventions}
Let $(C, d)$ be a differential graded object.
\begin{enumerate}[label=(\roman*)]
\item \emph{Homological grading.} $C = \bigoplus_{n \in \bZ} C_n$ with
$d: C_n \to C_{n-1}$. The differential \emph{lowers} degree by 1.
\item \emph{Cohomological grading.} $C = \bigoplus_{n \in \bZ} C^n$ with
$d: C^n \to C^{n+1}$. The differential \emph{raises} degree by 1.
\end{enumerate}
The translation is $C^n = C_{-n}$, which exchanges upper and lower indices.
\end{definition}

\begin{convention}[Standard conventions by context]
\label{conv:standard-by-context}
In this monograph:
\begin{enumerate}[label=(\roman*)]
\item \emph{Chain complexes of modules.} Homological convention $(C_\bullet, \partial)$.
\item \emph{Cochain complexes (de Rham, singular cochains).}
Cohomological convention $(C^\bullet, d)$.
\item \emph{Operadic bar/cobar.} Cohomological convention (translated from the homological convention of Loday--Vallette~\cite{LV12} via $C^n = C_{-n}$; see Definition~\ref{def:grading-conventions}).
\item \emph{Hochschild (co)homology.} Homological for homology $HH_*$,
cohomological for cohomology $HH^*$.
\item \emph{D-modules.} Cohomological convention, with de Rham complex
$\DR(\cM)$ in non-negative degrees.
\item \emph{$\infty$-categories.} Cohomological convention following
Lurie~\cite{HTT, HA}.
\end{enumerate}
\end{convention}

\begin{proposition}[Duality and grading reversal; \ClaimStatusProvedHere]
\label{prop:duality-grading}
For a finite-dimensional chain complex $(C_\bullet, \partial)$, the dual 
complex $C^* = \Hom(C, k)$ naturally carries a cohomological grading:
\[
(C^*)^n := \Hom_k(C_n, k) = (C_n)^*.
\]
The dual differential $\partial^*: (C^*)^n \to (C^*)^{n+1}$ is defined by:
\[
\langle \partial^* \phi, c \rangle := (-1)^{|\phi|+1} \langle \phi, \partial c \rangle
\]
for $\phi \in (C^*)^n$ and $c \in C_{n+1}$.
\end{proposition}

\begin{proof}
The sign ensures $(\partial^*)^2 = 0$. Indeed:
\begin{align*}
\langle (\partial^*)^2 \phi, c \rangle 
&= (-1)^{|\phi|+2} \langle \partial^* \phi, \partial c \rangle 
= (-1)^{|\phi|+2} (-1)^{|\phi|+1} \langle \phi, \partial^2 c \rangle = 0.
\end{align*}
The sign $(-1)^{|\phi|+1}$ is precisely the convention that makes
$(C, C^*)$ into a duality pairing of chain complexes.
\end{proof}

\begin{definition}[Connective and coconnective]
\label{def:connective}
A chain complex $C$ (in homological grading) is:
\begin{enumerate}[label=(\roman*)]
\item \emph{Connective} if $C_n = 0$ for $n < 0$ (concentrated in non-negative degrees).
\item \emph{Coconnective} if $C_n = 0$ for $n > 0$ (concentrated in non-positive degrees).
\item \emph{Bounded} if both conditions hold with finitely many nonzero terms.
\end{enumerate}
In cohomological grading, connective means $C^n = 0$ for $n > 0$, and 
coconnective means $C^n = 0$ for $n < 0$.
\end{definition}

\begin{remark}[Filtrations and spectral sequences]
\label{rem:filt-ss-grading}
The grading convention affects the indexing of spectral sequences. For a 
filtered complex with cohomological grading and increasing filtration 
$F^p C \subseteq F^{p+1} C$, the spectral sequence has:
\[
E_r^{p,q} \Rightarrow H^{p+q}(C).
\]
For homological grading with decreasing filtration $F_p C \supseteq F_{p+1} C$:
\[
E^r_{p,q} \Rightarrow H_{p+q}(C).
\]
The bidegree conventions $(p, q)$ versus $(p, n-p)$ also vary by author.
\end{remark}


\section{Suspensions and desuspensions}
\label{sec:susp-desusp}


\begin{definition}[Suspension and desuspension]
\label{def:suspension}
\index{suspension!conventions}
\index{desuspension!conventions}
For a graded object $V = \bigoplus_n V^n$:
\begin{enumerate}[label=(\roman*)]
\item The \emph{suspension} $sV$ (also denoted $V[-1]$ in the cohomological convention or $V[1]$ in the homological convention; see Remark~\ref{rem:shift-notation}) is
defined by:
\[
(sV)^n := V^{n-1}.
\]
An element $v \in V^{n-1}$ corresponds to $sv \in (sV)^n$ with $|sv| = |v| + 1$.

\item The \emph{desuspension} $s^{-1}V$ (also denoted $V[1]$ in the cohomological convention or $V[-1]$ in the homological convention) is defined by:
\[
(s^{-1}V)^n := V^{n+1}.
\]
An element $v \in V^{n+1}$ corresponds to $s^{-1}v \in (s^{-1}V)^n$ with 
$|s^{-1}v| = |v| - 1$.
\end{enumerate}
These are inverse operations: $s \circ s^{-1} = s^{-1} \circ s = \id$.

The two constructions that produce the Koszul dual use \emph{opposite} shifts, for a precise reason:
\begin{itemize}
\item \emph{Bar construction} (desuspension): $\bar{B}(\cA) = T^c(s^{-1}\bar{\cA},\, d)$.  Desuspension lowers generators by one degree because the bar differential $d_{\mathrm{res}}$ extracts a residue---an integration operation that lowers form degree---and the shift compensates so that $|d| = +1$.
\item \emph{Koszul dual coalgebra} (suspension): $\cA^! \hookrightarrow \mathrm{Cofree}(sV^*)$ (Appendix~\ref{app:koszul-reference}).  Suspension raises the dual generators $V^*$ by one degree because dualizing reverses the grading, and the shift restores it.
\end{itemize}
These are consistent: when $\cA$ is Koszul, $H^*(\bar{B}(\cA))$ concentrates in bar degree~$1$, and the desuspension-then-suspension cancels, recovering the Koszul dual generators in their natural degree.
\end{definition}

\begin{remark}[Shift notation ambiguity]
\label{rem:shift-notation}
The notation $V[n]$ is used with two incompatible conventions in the literature:
\begin{enumerate}[label=(\roman*)]
\item \emph{Homological convention.} $(V[n])_k = V_{k-n}$, so $V[1]$ shifts
\emph{up} (increases indices).
\item \emph{Cohomological convention.} $(V[n])^k = V^{k+n}$, so $V[1]$ shifts 
\emph{down} (decreases indices in the sense that degree-$k$ elements come from 
degree-$(k+1)$ elements of $V$).
\end{enumerate}
We follow the cohomological convention: $V[n]^k = V^{k+n}$. Under this convention,
$sV = V[-1]$ (since $(V[-1])^k = V^{k-1} = (sV)^k$) and
$s^{-1}V = V[1]$ (since $(V[1])^k = V^{k+1} = (s^{-1}V)^k$).

\emph{Warning.} Many references in algebraic topology use the opposite convention where $sV = V[1]$; this corresponds to $(V[n])^k = V^{k-n}$ (the homological convention). When importing results from such references, all shift indices must be negated.
\end{remark}

\begin{proposition}[Suspension and differentials; \ClaimStatusProvedHere]
\label{prop:susp-diff}
If $(V, d)$ is a chain complex with $d: V^n \to V^{n+1}$, then $sV = V[-1]$
carries the \emph{suspended differential}:
\[
d_{sV}: (sV)^n \to (sV)^{n+1}, \quad d_{sV}(sv) := -s(dv).
\]
The sign ensures the suspension defines a functor on chain complexes (compatible with tensor products and the Koszul sign rule).
\end{proposition}

\begin{proof}
Without the sign, we would have $d_{sV}(sv) = s(dv)$, and then:
\[
d_{sV}^2(sv) = d_{sV}(s(dv)) = s(d(dv)) = s(d^2 v) = 0.
\]
This appears to work, but the sign is required for compatibility with tensor 
products. Consider $V \otimes W$ with differential:
\[
d_{V \otimes W}(v \otimes w) = dv \otimes w + (-1)^{|v|} v \otimes dw.
\]
For the suspended tensor product $(sV) \otimes W$ to have consistent differential:
\[
d((sv) \otimes w) = d_{sV}(sv) \otimes w + (-1)^{|sv|} sv \otimes dw 
= -s(dv) \otimes w + (-1)^{|v|+1} sv \otimes dw.
\]
This matches $s(d(v \otimes w))$ only with the sign convention $d_{sV}(sv) = -s(dv)$.
\end{proof}

\begin{corollary}[Iterated suspension \cite{LV12}; \ClaimStatusProvedElsewhere]
\label{cor:iterated-susp}
For $n$-fold suspension $s^n V = V[-n]$:
\[
d_{s^n V}(s^n v) = (-1)^n s^n(dv).
\]
The suspended element $s^n v$ has degree $|s^n v| = |v| + n$.
\end{corollary}

\begin{definition}[Suspension in operadic context]
\label{def:susp-operadic}
For an operad $\cP$, the \emph{operadic suspension} $\fS \cP$ is the operad with:
\[
(\fS \cP)(n) := \sgn_n \otimes \cP(n)[n-1]
\]
where $\sgn_n$ is the sign representation of $\Sigma_n$. The operadic composition 
involves signs from both the degree shift and the sign representation.

Explicitly, if $\gamma \in \cP(k)$ and $\mu_1, \ldots, \mu_k \in \cP(n_1), \ldots, \cP(n_k)$:
\[
s^{k-1}\gamma \circ (s^{n_1-1}\mu_1, \ldots, s^{n_k-1}\mu_k) 
= \epsilon \cdot s^{n-1}(\gamma \circ (\mu_1, \ldots, \mu_k))
\]
where $n = n_1 + \cdots + n_k$ and $\epsilon$ is a sign depending on degrees 
and positions.
\end{definition}

\begin{proposition}[Suspension and Koszul duality \cite{LV12}; \ClaimStatusProvedElsewhere]
\label{prop:susp-koszul}
For a Koszul operad $\cP$, the Koszul dual cooperad is:
\[
\cP^{\scriptstyle \text{\textexclamdown}} = (\fS^{-1} \cP^!)^\vee
\]
where $\cP^!$ is the linear dual operad and $\fS^{-1}$ is operadic desuspension. 
Explicitly:
\[
\cP^{\scriptstyle \text{\textexclamdown}}(n) = \cP^!(n)^* \otimes \sgn_n[1-n].
\]
This formula explains the appearance of signs and shifts in bar-cobar duality.
\end{proposition}


\section{Determinant lines}
\label{sec:det-lines}

\begin{definition}[Determinant of a vector space]
\label{def:det-vect}
For a finite-dimensional vector space $V$ of dimension $n$, the 
\emph{determinant line} is:
\[
\Det(V) := \bigwedge^n V = \bigwedge^{\dim V} V.
\]
This is a 1-dimensional vector space. An element $\omega \in \Det(V) \setminus \{0\}$ 
is called an \emph{orientation} of $V$.
\end{definition}

\begin{proposition}[Properties of determinant lines \cite{Weibel94}; \ClaimStatusProvedElsewhere]
\label{prop:det-properties}
The determinant construction satisfies:
\begin{enumerate}[label=(\roman*)]
\item \emph{Short exact sequences.} For $0 \to U \to V \to W \to 0$ exact,
\[
\Det(V) \cong \Det(U) \otimes \Det(W).
\]
\item \emph{Direct sums.} $\Det(V \oplus W) \cong \Det(V) \otimes \Det(W)$.
\item \emph{Duality.} $\Det(V^*) \cong \Det(V)^* = \Det(V)^{-1}$.
\item \emph{Zero space.} $\Det(0) = k$ (the ground field, concentrated in degree 0).
\item \emph{Functoriality.} An isomorphism $f: V \xrightarrow{\sim} W$ induces 
$\Det(f): \Det(V) \xrightarrow{\sim} \Det(W)$ with $\Det(f)(v_1 \wedge \cdots \wedge v_n) 
= f(v_1) \wedge \cdots \wedge f(v_n)$.
\end{enumerate}
\end{proposition}

\begin{definition}[Determinant of a finite set]
\label{def:det-set}
For a finite set $I$, define:
\[
\Det(I) := \Det(k^I)
\]
where $k^I$ is the vector space with basis indexed by $I$. Concretely, if 
$I = \{i_1, \ldots, i_n\}$:
\[
\Det(I) = k \cdot (e_{i_1} \wedge \cdots \wedge e_{i_n})
\]
is 1-dimensional, and permuting the order of $I$ introduces signs according to 
the permutation's signature.
\end{definition}

\begin{lemma}[Determinant and ordering \cite{Weibel94}; \ClaimStatusProvedElsewhere]
\label{lem:det-ordering}
For a finite set $I$ with total orderings $<$ and $<'$, let $\sigma \in \Sigma_I$ 
be the permutation taking the $<$-ordering to the $<'$-ordering. Then:
\[
e_{i_1} \wedge \cdots \wedge e_{i_n} = \sgn(\sigma) \cdot e_{i_{\sigma(1)}} \wedge \cdots \wedge e_{i_{\sigma(n)}}.
\]
The determinant line $\Det(I)$ is canonically independent of ordering, but a 
choice of generator (i.e., orientation) depends on ordering.
\end{lemma}

\begin{definition}[Orientation line]
\label{def:orientation-line}
For a smooth manifold $M$ of dimension $d$, the \emph{orientation line} at 
$p \in M$ is:
\[
\orline{M,p} := \Det(T_p M) = \bigwedge^d T_p M.
\]
The \emph{orientation sheaf} $\orline{M}$ is the local system with fiber 
$\orline{M,p}$ at $p$. An orientation of $M$ is a global section 
$\sigma: M \to \orline{M}$ that is nowhere vanishing.
\end{definition}

\begin{proposition}[Determinant lines on configuration spaces \cite{FM94}; \ClaimStatusProvedElsewhere]
\label{prop:det-config}
For the configuration space $\Conf_n(X)$ of a curve $X$:
\begin{enumerate}[label=(\roman*)]
\item The tangent space at $(z_1, \ldots, z_n) \in \Conf_n(X)$ is:
\[
T_{(z_1, \ldots, z_n)} \Conf_n(X) = \bigoplus_{i=1}^n T_{z_i} X.
\]
\item The determinant line is:
\[
\Det(T\Conf_n(X)) = \bigotimes_{i=1}^n \Det(T_{z_i} X) = \bigotimes_{i=1}^n \omega_{X,z_i}^{-1}
\]
where $\omega_X = \Det(T^*X)$ is the canonical bundle.
\item For the Fulton--MacPherson compactification $\FM_n(X)$:
\[
\Det(T\FM_n(X)) \cong \Det(T\Conf_n(X)) \otimes \Det(N_{D/\FM_n})
\]
where $N_{D/\FM_n}$ is the normal bundle to the boundary divisor.
\end{enumerate}
\end{proposition}

\begin{construction}[Determinant line in bar complex]
\label{constr:det-bar}
In the geometric bar complex, the determinant line appears through the 
identification:
\[
\Bbar^{\mathrm{geom}}_n(\cA) = \Gamma(\FM_n(X), \cA^{\boxtimes n} \otimes \Omega^{n-1}_{\log}(\FM_n, D))
\]
The logarithmic $(n-1)$-forms can be written as:
\[
\Omega^{n-1}_{\log}(\FM_n, D) \cong \omega_{\FM_n} \otimes \Det(N_D)^{-1}[1-n]
\]
tracking the relationship between the canonical bundle and boundary geometry.

Explicitly, for $n = 2$ on $X = \bP^1$:
\[
\Omega^1_{\log}(\FM_2(\bP^1), D_{12}) = \cO \cdot \dlog(z_1 - z_2) = \cO \cdot \frac{d(z_1 - z_2)}{z_1 - z_2}
\]
which has a simple pole along $D_{12} = \{z_1 = z_2\}$ and generates the 
log de Rham complex.
\end{construction}

\begin{proposition}[Determinant and residues \cite{Har77}; \ClaimStatusProvedElsewhere]
\label{prop:det-residue}
The residue map along a divisor $D \subset Y$ is:
\[
\Res_D: \Omega^k_{\log}(Y, D) \to \Omega^{k-1}(D)
\]
This map respects determinant lines in the following sense: if $D$ is smooth 
of codimension 1, then:
\[
\Res_D: \omega_Y \otimes \cO_Y(D)|_D \to \omega_D
\]
is the adjunction isomorphism, and the signs in the bar differential arise 
from tracking this isomorphism across multiple boundary strata.
\end{proposition}

\begin{theorem}[Determinant conventions and bar-cobar signs \cite{LV12}; \ClaimStatusProvedElsewhere]
\label{thm:det-bar-cobar-signs}
The signs in the bar-cobar adjunction can be uniformly expressed using 
determinant lines. For an $\Eone$-chiral algebra $\cA$:
\begin{enumerate}[label=(\roman*)]
\item The bar complex element $[a_1 | \cdots | a_n] \in \Bbar_n(\cA)$ corresponds to:
\[
a_1 \otimes \cdots \otimes a_n \otimes \omega_{[n]} \in \cA^{\otimes n} \otimes \Det([n])
\]
where $[n] = \{1, \ldots, n\}$ and $\omega_{[n]}$ is the standard generator of 
$\Det([n])$ determined by the ordering $1 < 2 < \cdots < n$.

\item The bar differential is:
\[
d[a_1 | \cdots | a_n] = \sum_{i=1}^{n-1} (-1)^{\epsilon_i} [a_1 | \cdots | a_i a_{i+1} | \cdots | a_n]
\]
where $\epsilon_i = |a_1| + \cdots + |a_i| + i - 1$ tracks both internal degrees 
and position.

\item These signs arise canonically from the face maps of the simplicial 
structure on $\Delta^\bullet$ and the determinant line conventions.
\end{enumerate}
\end{theorem}


\section{Detailed sign computations in bar-cobar duality}
\label{sec:detailed-signs}

\begin{computation}[Bar differential signs: degree 2 to 1]
\label{comp:bar-signs-2-1}
For elements $[a|b] \in \Bbar_2(A)$ where $A$ is a dg-algebra:
\[
d_{\mathrm{bar}}[a|b] = [ab] - (-1)^{|a|} a \cdot [b] - [a] \cdot b
\]
Here:
\begin{enumerate}[label=(\roman*)]
\item The term $[ab]$ comes from the multiplication $\mu: A \otimes A \to A$.
\item The term $(-1)^{|a|} a \cdot [b]$ is the left action of $A$ on $\Bbar_1(A) = A$.
\item The term $[a] \cdot b$ is the right action.
\end{enumerate}

For the \emph{reduced} bar complex $\overline{\Bbar}(A) = \Bbar(A, A, k)$:
\[
d_{\mathrm{bar}}[a|b] = [ab]
\]
with augmentation eliminating the action terms.
\end{computation}

\begin{computation}[Bar differential signs: degree 3 to 2]
\label{comp:bar-signs-3-2}
For $[a|b|c] \in \Bbar_3(A)$:
\begin{align*}
d_{\mathrm{bar}}[a|b|c] &= [ab|c] - (-1)^{|a|}[a|bc] \\
&\quad + (-1)^{|a|} a \cdot [b|c] + (-1)^{|a|+|b|} [a|b] \cdot c \\
&\quad + (-1)^{|a|} [da|b|c] + (-1)^{|a|+|b|+1} [a|db|c] + (-1)^{|a|+|b|+|c|} [a|b|dc]
\end{align*}

The signs arise systematically:
\begin{enumerate}[label=(\roman*)]
\item $[ab|c]$: Multiply positions 1,2; no sign (conventions place multiplication first).
\item $(-1)^{|a|}[a|bc]$: Multiply positions 2,3; sign from passing $a$ over boundary.
\item Action terms: Signs from Koszul rule and boundary conventions.
\item Differential terms: Signs from position in the bar expression.
\end{enumerate}
\end{computation}

\begin{computation}[Cobar differential signs: degree 1 to 2]
\label{comp:cobar-signs-1-2}
For a coassociative coalgebra $C$ with coproduct $\Delta$, the cobar construction 
$\Omega(C) = T(s^{-1}\overline{C})$ has differential on generators:
\[
d_\Omega(s^{-1}c) = -s^{-1}(dc) + \sum_{(c)} (-1)^{|c'|} (s^{-1}c') \otimes (s^{-1}c'')
\]
where $\Delta(c) = \sum_{(c)} c' \otimes c''$ (Sweedler notation) and the desuspended element $s^{-1}c'$ has degree $|s^{-1}c'| = |c'| - 1$.

\emph{Sign verification.} The term $(-1)^{|c'|}$ arises from:
\[
d_\Omega(s^{-1}c) \propto (s^{-1} \otimes s^{-1}) \circ \Delta(c) = \sum (-1)^{|s^{-1}| \cdot |c'|} s^{-1}c' \otimes s^{-1}c''
\]
Since $|s^{-1}| = -1$, we get $(-1)^{-|c'|} = (-1)^{|c'|}$ (as $(-1)^{-n} = (-1)^n$).
\end{computation}

\begin{computation}[Cobar differential signs: degree 2 to 3]
\label{comp:cobar-signs-2-3}
For a tensor $s^{-1}c_1 \otimes s^{-1}c_2 \in \Omega_2(C)$:
\begin{align*}
d_\Omega(s^{-1}c_1 \otimes s^{-1}c_2) &= d_\Omega(s^{-1}c_1) \otimes s^{-1}c_2 
    + (-1)^{|s^{-1}c_1|} s^{-1}c_1 \otimes d_\Omega(s^{-1}c_2) \\
&= \Bigl( -s^{-1}(dc_1) + \sum (-1)^{|c_1'|} s^{-1}c_1' \otimes s^{-1}c_1'' \Bigr) \otimes s^{-1}c_2 \\
&\quad + (-1)^{|c_1|-1} s^{-1}c_1 \otimes \Bigl( -s^{-1}(dc_2) + \sum (-1)^{|c_2'|} s^{-1}c_2' \otimes s^{-1}c_2'' \Bigr)
\end{align*}

Expanding:
\begin{align*}
d_\Omega(s^{-1}c_1 \otimes s^{-1}c_2) &= -s^{-1}(dc_1) \otimes s^{-1}c_2 \\
&\quad + \sum (-1)^{|c_1'|} s^{-1}c_1' \otimes s^{-1}c_1'' \otimes s^{-1}c_2 \\
&\quad + (-1)^{|c_1|} s^{-1}c_1 \otimes s^{-1}(dc_2) \\
&\quad + (-1)^{|c_1|-1} \sum (-1)^{|c_2'|} s^{-1}c_1 \otimes s^{-1}c_2' \otimes s^{-1}c_2''
\end{align*}
\end{computation}

\begin{proposition}[Master sign formula {\cite{LV12}}; \ClaimStatusProvedElsewhere]
\label{prop:master-sign}
For the bar complex of an $\Ainf$-algebra $(A, \{m_n\})$, the full differential 
on $[a_1|\cdots|a_n]$ is:
\[
d_{\mathrm{bar}}[a_1|\cdots|a_n] = \sum_{i=1}^{n} \sum_{k=1}^{n-i+1} 
(-1)^{\epsilon_{i,k}} [a_1|\cdots|a_{i-1}|m_k(a_i, \ldots, a_{i+k-1})|a_{i+k}|\cdots|a_n]
\]
where the sign is:
\[
\epsilon_{i,k} = \sum_{j=1}^{i-1} (|a_j| - 1) + \sum_{j=i}^{i+k-2} |a_j|
\]
where the first sum is the \emph{coderivation sign} (from commuting past
$s^{-1}a_1, \ldots, s^{-1}a_{i-1}$) and the second is the \emph{desuspension sign}
from $b_k(s^{-1}a_i, \ldots, s^{-1}a_{i+k-1})$.

\begin{remark}[Sign convention comparison]
Some references (e.g., Loday--Vallette \cite{LV12}) define the bar differential
on $sa_1 \otimes \cdots \otimes sa_n$ with a sign
$\epsilon'_{i,k} = \sum_{j=1}^{i-1}|sa_j|$ arising from commuting the coderivation past
suspended elements.  For a strict algebra ($m_k = 0$ for $k \neq 2$), the two conventions
are related by a global sign on each summand and yield the same $d^2 = 0$ complex.
The explicit computations in \S\ref{comp:bar-signs-2-1}--\ref{comp:bar-signs-3-2}
use the \emph{unsuspended} bar notation ($[a|b] = a \otimes b$).  In the desuspended convention ($[a|b] = s^{-1}a \otimes s^{-1}b$), the correct formulas are
$d[a|b] = (-1)^{|a|}[ab]$ and $d[a|b|c] = (-1)^{|a|}[ab|c] - (-1)^{|a|+|b|}[a|bc]$; see the verification in \S\ref{ver:d2-degree3}.
\end{remark}
\end{proposition}

\medskip
\noindent\textbf{Convention shift.}  The computations above (\S\ref{comp:bar-signs-2-1}--\ref{comp:bar-signs-3-2}) use the unsuspended bar notation $[a|b] = a \otimes b$.  From this point onward, we work in the \emph{desuspended} convention $[a|b] = s^{-1}a \otimes s^{-1}b$, which is the convention used in the body of the monograph.

\begin{verification}[\texorpdfstring{$d^2 = 0$}{d\textasciicircum 2 = 0} check at degree 3]
\label{ver:d2-degree3}
We verify $d_{\mathrm{bar}}^2[a|b|c] = 0$ for a graded associative algebra
(where $m_2 = \mu$ and $m_k = 0$ for $k \neq 2$).

The desuspended multiplication on $s^{-1}A$ carries a Koszul sign:
$b_2(s^{-1}a, s^{-1}b) = (-1)^{|a|} s^{-1}(ab)$.  The coderivation extension
to $T^c(s^{-1}\bar{A})$ gives:
\begin{align*}
d_{\mathrm{bar}}[a|b] &= (-1)^{|a|}[ab], \\
d_{\mathrm{bar}}[a|b|c] &= (-1)^{|a|}[ab|c] - (-1)^{|a|+|b|}[a|bc].
\end{align*}
(The sign at position $i$ is $(-1)^{\sum_{j<i}(|a_j|-1) + |a_i|}$.)

Applying $d_{\mathrm{bar}}$ again:
\begin{align*}
d_{\mathrm{bar}}^2[a|b|c] &= (-1)^{|a|} d_{\mathrm{bar}}[ab|c]
    - (-1)^{|a|+|b|} d_{\mathrm{bar}}[a|bc] \\
&= (-1)^{|a|} \cdot (-1)^{|a|+|b|} [(ab)c]
    - (-1)^{|a|+|b|} \cdot (-1)^{|a|} [a(bc)] \\
&= (-1)^{|b|}\bigl[(ab)c - a(bc)\bigr] = 0
\end{align*}
by associativity.

\begin{remark}
For generators concentrated in degree~0 (the ungraded case), all signs are
trivial and the formulas simplify to $d[a|b] = [ab]$ and
$d[a|b|c] = [ab|c] - [a|bc]$, recovering the expressions in
Computations~\ref{comp:bar-signs-2-1}--\ref{comp:bar-signs-3-2}.
\end{remark}
\end{verification}


\section{Determinant lines and stratifications}
\label{sec:det-strat}

\begin{definition}[Orientation system]
\label{def:orientation-system}
An \emph{orientation system} on a stratified space $Y = \bigsqcup_\alpha Y_\alpha$ 
is a collection of line bundles $\orline{Y_\alpha}$ on each stratum together with 
\emph{gluing isomorphisms}: for each pair $Y_\alpha \subseteq \overline{Y_\beta}$ 
(closure relation), an isomorphism:
\[
\phi_{\alpha\beta}: \orline{Y_\alpha}|_{\partial_\alpha Y_\beta} \xrightarrow{\sim} 
\orline{Y_\beta}|_{\partial_\alpha Y_\beta} \otimes N_{\alpha\beta}
\]
where $N_{\alpha\beta}$ is the normal bundle factor and $\partial_\alpha Y_\beta$ 
is the boundary of $Y_\beta$ meeting $Y_\alpha$.
\end{definition}

\begin{proposition}[Orientation system on FM compactification \cite{FM94}; \ClaimStatusProvedElsewhere]
\label{prop:orient-fm}
The FM compactification $\FM_n(X)$ carries a canonical orientation system 
determined by:
\begin{enumerate}[label=(\roman*)]
\item On the open stratum $\Conf_n(X)$: $\orline{\Conf_n(X)} = \bigotimes_{i=1}^n \orline{X}$.
\item On boundary stratum $D_T$ (labeled by tree $T$): 
$\orline{D_T} = \bigotimes_{v \in V(T)} \orline{X}$ times exceptional fiber orientations.
\item Gluing isomorphisms: Given by the blowup construction, with signs determined 
by the ordering on vertices of $T$.
\end{enumerate}
\end{proposition}

\begin{construction}[Sign from blowup]
\label{constr:sign-blowup}
When blowing up the diagonal $\Delta_{ij} \subset X^n$, the orientation 
transformation is:
\[
\orline{X^n} = \orline{\Bl_{\Delta_{ij}}(X^n)} \otimes \orline{E_{ij}}^{-1}
\]
where $E_{ij} \cong \bP(N_{\Delta_{ij}/X^n})$ is the exceptional divisor.

For $X$ a curve (dimension 1), the normal bundle $N_{\Delta/X^2} \cong T_X$ 
along the diagonal, so:
\[
\orline{E_{ij}} = \orline{\bP^0} = k
\]
is canonically trivial, and no orientation sign arises from the blowup in this case.
\end{construction}

\begin{lemma}[Residue and orientation \cite{FM94, Har77}; \ClaimStatusProvedElsewhere]
\label{lem:residue-orient}
The residue map $\Res_{D_{ij}}: \Omega^k_{\log}(\FM_n, D) \to \Omega^{k-1}(D_{ij})$
is compatible with orientations:
\[
\Res_{D_{ij}}(\omega \wedge \dlog(z_i - z_j)) = \omega|_{D_{ij}}
\]
with sign conventions:
\begin{enumerate}[label=(\roman*)]
\item If $\omega = f(z_1, \ldots, z_n) \, dz_1 \wedge \cdots \wedge \widehat{dz_i} \wedge \cdots \wedge dz_n$
(omitting $dz_i$), then:
\[
\Res_{D_{ij}}(\omega \wedge \dlog(z_i - z_j)) = (-1)^{i-1} f|_{z_i = z_j} \,
dz_1 \wedge \cdots \wedge \widehat{dz_i} \wedge \cdots \wedge dz_n|_{D_{ij}}.
\]
\item The sign $(-1)^{i-1}$ accounts for moving $\dlog(z_i - z_j)$ to the position
where $dz_i$ would appear.
\end{enumerate}
\end{lemma}


%% ======================================================================
%% Dictionary of sign conventions
%% (Merged from former appendices/sign_conventions.tex)
%% ======================================================================

\section{Dictionary of sign conventions}
\label{app:sign-conventions}

Sign conventions used in different sources are compared and translated below.

\subsection{Loday--Vallette vs. this manuscript}

\begin{center}
\begin{longtable}{|>{\raggedright\arraybackslash}p{5.8cm}|>{\raggedright\arraybackslash}p{3.9cm}|>{\raggedright\arraybackslash}p{3.9cm}|}
\caption{Sign convention comparison} \\
\hline
\textbf{Object/Operation} & \textbf{Loday--Vallette \cite{LV12}} & \textbf{This Manuscript} \\
\hline
\endfirsthead
\multicolumn{3}{c}{\tablename\ \thetable\ -- \textit{Continued}} \\
\hline
\textbf{Object/Operation} & \textbf{Loday--Vallette} & \textbf{This Manuscript} \\
\hline
\endhead
\hline
\endfoot
\hline
\endlastfoot
Koszul sign rule & $(-1)^{|a| \cdot |b|}$ & $(-1)^{|a| \cdot |b|}$ (same) \\
Differential degree & $|d| = -1$ & $|d| = +1$ \\
Suspension & $s: V \to sV$, $|sv| = |v| + 1$ & $s: V \to V[-1]$, $|sv| = |v| + 1$ (cohomological: $V[-1]^k = V^{k-1}$, so $v \in V^m$ maps to $sV^{m+1}$) \\
Desuspension & $s^{-1}: sV \to V$, $|s^{-1}(sv)| = |v|$ & $[-1]: V[1] \to V$ \\
Bar construction & $B(\mathcal{A}) = (T^c(s\mathcal{A}), d_{\text{bar}})$ & $\bar{B}(\mathcal{A}) = (T^c(s^{-1}\bar{\mathcal{A}}), d)$ (desuspension; cf.\ cohomological convention) \\
Cobar construction & $\Omega(\mathcal{C}) = (T(s^{-1}\mathcal{C}), d_{\text{cobar}})$ & $\Omega(\mathcal{C})_n = \mathcal{C}^{\otimes_{\text{co}} n}$ \\
Coproduct sign & $\Delta(ab) = \Delta(a) \cdot \Delta(b)$ with $(-1)^{|a_2| \cdot |b_1|}$ & Same \\
Shuffle sign & $(a \shuffle b) = \sum (-1)^{\sigma} a_{\sigma(1)} \otimes \cdots$ & Same \\
Collision divisor ordering & Not applicable (no geometry) & Lexicographic: $i < j$ always \\
Logarithmic form sign & Not applicable & $\eta_{ji} = \eta_{ij}$ as differential forms ($d\log(-1)=0$). Note: the Orlik--Solomon generators are antisymmetric; fixing $i < j$ in $\eta_{ij}$ fixes an orientation. \\
Arnold relations & Not applicable & $\sum_{\text{cyclic}} \eta_{ij} \wedge \eta_{jk} = 0$ \\
\end{longtable}
\end{center}

\subsubsection{Key differences explained}

\emph{1.~Differential degree.}
Loday--Vallette use homological grading ($|d| = -1$); this manuscript uses cohomological grading ($|d| = +1$).

\emph{Translation.} If $(C_*, d)$ is a chain complex in LV convention ($d: C_n \to C_{n-1}$), set $C^n := C_{-n}$ with the same underlying map~$d$, which now raises cohomological degree by~$1$.  No sign change is needed: $d_{\text{us}} = d_{\text{LV}}$ (as linear maps).

\emph{2.~Suspension.}
Loday--Vallette treat suspension $s$ as an explicit operator; this manuscript uses the degree shift $sV = V[-1]$, so $|sv| = |v| + 1$.

\emph{Translation.} $sV$ (LV) corresponds to $V[-1]$ (us), i.e., our $sV$.  Both raise degree by~$1$: LV's convention $(sV)_k = V_{k-1}$ translates via $\deg_{\mathrm{coh}} = -\deg_{\mathrm{hom}}$ to $(sV)^k = V^{k-1} = (V[-1])^k$.

\emph{3.~Bar/cobar formulas.}
Loday--Vallette include explicit suspension in the formula; this manuscript absorbs it into the D-module structure.

\emph{Translation.} Our $\bar{B}_n(\mathcal{A})$ corresponds to LV's $B(\mathcal{A})$ in degree $n$ after removing the suspension $s$.

\subsection{Beilinson--Drinfeld vs. this manuscript}

\begin{center}
\begin{tabular}{|>{\raggedright\arraybackslash}p{5.8cm}|>{\raggedright\arraybackslash}p{3.9cm}|>{\raggedright\arraybackslash}p{3.9cm}|}
\hline
\textbf{Object/Operation} & \textbf{BD \cite{BD04}} & \textbf{This Manuscript} \\
\hline
Chiral algebra & $\mathcal{A}$ (D-module) & $\mathcal{A}$ (same) \\
Factorization & $j_{!*}\mathcal{A}$ & Implicit in $\ConfigSpace{n}$ \\
Configuration space & $\text{Ran}(X)$ & $\bigcup_n \ConfigSpace{n}$ \\
Collision divisor & $(i,j)^c$ (complement) & $\omitpair{i}{j}$ (hat notation) \\
Chiral connection & $\nabla$ & $\nabla$ (same) \\
Residue & Implicit in factorization & Explicit via $\text{Res}_{z_i \to z_j}$ \\
OPE & Encoded in $\mathcal{A}(D)$ & Explicit: $a(z)b(w) = \sum c_n(w)/(z-w)^n$ \\
Koszul duality & Not discussed (BD focus on one side) & Central theme \\
\hline
\end{tabular}
\end{center}

\subsubsection{Key differences explained}

\emph{BD's approach.}
\begin{itemize}
\item Ran space formulation (taking colimits over all $n$)
\item Factorization axiom as the primary structure
\item D-modules as the fundamental objects
\item No explicit bar/cobar constructions
\end{itemize}

\emph{This monograph.}
\begin{itemize}
\item Explicit configuration spaces $\ConfigSpace{n}$ for each $n$
\item Bar complex as an explicit resolution
\item Koszul duality as the organizing principle
\item Geometric realization via residues
\end{itemize}

\emph{Translation.} Our bar complex $\bar{B}_*(\mathcal{A})$ is the ``Chevalley--Cousin resolution'' of $\mathcal{A}$ viewed as a factorization algebra in BD's sense (BD Theorem 3.4.9).

\subsection{Costello--Gwilliam vs. this manuscript}

\begin{center}
\begin{tabular}{|>{\raggedright\arraybackslash}p{5.8cm}|>{\raggedright\arraybackslash}p{3.9cm}|>{\raggedright\arraybackslash}p{3.9cm}|}
\hline
\textbf{Object/Operation} & \textbf{CG \cite{CG17}} & \textbf{This Manuscript} \\
\hline
Factorization algebra & $\mathcal{F}$ (precosheaf) & $\mathcal{A}$ (chiral algebra) \\
Configuration space & $\text{Conf}_n(M)$ & $\ConfigSpace{n}$ \\
Compactification & Fulton--MacPherson & Same \\
Bar complex & Not emphasized & Central construction \\
Feynman diagrams & Graph complex & Genus-stratified bar complex \\
Renormalization & BV formalism & Completion of bar complex \\
Quantum corrections & Loop order & Genus $g$ \\
\hline
\end{tabular}
\end{center}

\subsubsection{Key differences explained}

\emph{CG's approach.}
\begin{itemize}
\item Factorization algebras on manifolds (any dimension)
\item Feynman diagram expansion for observables
\item BV formalism for quantization
\item Emphasis on renormalization and effective field theory
\end{itemize}

\emph{This monograph.}
\begin{itemize}
\item Chiral algebras on curves (1-dimensional)
\item Genus expansion instead of loop expansion
\item Bar-cobar construction and Koszul duality instead of BV
\item Emphasis on algebraic structure and representation theory
\end{itemize}

\emph{Translation.}
\begin{itemize}
\item CG's ``factorization algebra'' on a curve $X$ = BD/Our ``chiral algebra''
\item CG's ``observable'' = Our ``element of $\mathcal{A}$''
\item CG's ``1-loop graph'' = Our ``genus-1 contribution to bar complex''
\item CG's ``renormalization'' = Our ``nilpotent completion of bar complex''
\end{itemize}

\subsection{Kontsevich vs. this manuscript}

\begin{center}
\begin{tabular}{|>{\raggedright\arraybackslash}p{5.8cm}|>{\raggedright\arraybackslash}p{3.9cm}|>{\raggedright\arraybackslash}p{3.9cm}|}
\hline
\textbf{Object/Operation} & \textbf{Kontsevich \cite{Kon99}} & \textbf{This Manuscript} \\
\hline
Configuration space & $C_n(\mathbb{R}^d)$ (open) & $\ConfigSpace{n}$ (compactified) \\
Forms & Smooth forms & Logarithmic forms \\
Propagator & $\frac{1}{z-w}$ & $\frac{dz}{z-w}$ (with log pole) \\
Graphs & Directed graphs & Genus-stratified graphs \\
Formality & $\mathcal{U}: T_{\text{poly}} \to D_{\text{poly}}$ & Bar-cobar: $\bar{B} \dashv \Omega$ \\
Wheels & Non-planar graphs & Higher genus contributions \\
\hline
\end{tabular}
\end{center}

\subsubsection{Key differences explained}

\emph{Kontsevich's formality.}
\begin{itemize}
\item Classical: Poisson manifolds $\to$ deformation quantization
\item Configuration spaces in $\mathbb{R}^d$ (no compactification)
\item Smooth forms (no poles)
\item Graph complex (combinatorial)
\end{itemize}

\emph{Our chiral formality.}
\begin{itemize}
\item Coisson algebras (PVA) $\to$ vertex algebras (OPE) via chiral configuration spaces
\item Configuration spaces on curves (compactified)
\item Logarithmic forms (poles at collisions)
\item Genus-stratified complex (geometric)
\end{itemize}

\emph{Relation.} Kontsevich's formality is the \emph{genus-0 part} of chiral formality. The bar-cobar framework extends Kontsevich to include all genera (higher-loop corrections).

\subsection{Summary table: all conventions}

\begin{center}
\begin{longtable}{|>{\raggedright\arraybackslash}p{5cm}|>{\raggedright\arraybackslash}p{2cm}|>{\raggedright\arraybackslash}p{2cm}|>{\raggedright\arraybackslash}p{2cm}|>{\raggedright\arraybackslash}p{2cm}|}
\caption{Sign convention dictionary} \\
\hline
\textbf{Operation} & \textbf{LV} & \textbf{BD} & \textbf{CG} & \textbf{Us} \\
\hline
\endfirsthead
\multicolumn{5}{c}{\tablename\ \thetable\ -- \textit{Continued}} \\
\hline
\textbf{Operation} & \textbf{LV} & \textbf{BD} & \textbf{CG} & \textbf{Us} \\
\hline
\endhead
Koszul sign & $(-1)^{|a||b|}$ & $(-1)^{|a||b|}$ & $(-1)^{|a||b|}$ & $(-1)^{|a||b|}$ \\
Differential degree & $-1$ & $+1$ & $+1$ & $+1$ \\
Suspension degree & $+1$ & $+1$ & $+1$ & $+1$ \\
Collision sign & N/A & Implicit & Explicit & Explicit \\
Arnold relations & N/A & Implicit & Yes & Yes \\
\hline
\end{longtable}
\end{center}

\emph{Recommendation.} When reading other sources, always consult this appendix to translate signs and conventions. The most common source of error in chiral algebra computations is inconsistent sign conventions.

\subsection{Practical guide: how to translate}

\subsubsection{From Loday--Vallette to our conventions}

\emph{Rule~1.} Replace $|d| = -1$ with $|d| = +1$.

\emph{Rule~2.} Replace LV's suspension $s$ with our degree shift $[-1]$: $sV = V[-1]$.

\emph{Rule~3.} Our geometric forms $\Omega^n(\log D)$ correspond to LV's suspended generators $s^n V$.

\subsubsection{From Beilinson--Drinfeld to our conventions}

\emph{Rule~1.} BD's $j_{!*}$ (factorization pushforward) = our explicit configuration space integral.

\emph{Rule~2.} BD's Ran space $\text{Ran}(X)$ = our $\bigcup_n \ConfigSpace{n} / \sim$ (colimit).

\emph{Rule~3.} BD's ``chiral product'' = our ``OPE residue''

\subsubsection{From Costello--Gwilliam to our conventions}

\emph{Rule~1.} CG's ``1-loop'' = our ``genus 1''

\emph{Rule~2.} CG's ``renormalization'' = our ``$I$-adic completion''

\emph{Rule~3.} CG's BV bracket = our bar differential (when specialized to curves).

\subsection{Examples of translation}

\begin{example}[Translating a bar differential formula from LV]
\emph{Loday--Vallette writes.}
\[d_{LV}(s a_1 \otimes s a_2 \otimes s a_3) = (-1)^{|a_1|} (s[a_1, a_2] \otimes s a_3) + \cdots\]

\emph{Our notation.}
\[d(a_1 \otimes a_2 \otimes a_3 \otimes \omega_2) = ([a_1, a_2] \otimes a_3 \otimes \text{Res}[\omega_2]) + \cdots\]

\emph{Translation.}
\begin{itemize}
\item Remove explicit suspension $s$ (implicit in form degree)
\item Change sign convention: $(-1)^{|a_1|}$ becomes automatic from Koszul rule
\item Add geometric form $\omega_2 \in \Omega^2(\log D)$
\end{itemize}
\end{example}

\begin{example}[Translating a BD factorization formula]
\emph{Beilinson--Drinfeld writes.}
\[j_{!*}(\mathcal{A} \boxtimes \mathcal{A})|_{U \sqcup V} \simeq j_{!*}\mathcal{A}|_U \boxtimes j_{!*}\mathcal{A}|_V\]

\emph{Our notation.}
\[\barcomplex{1}{\mathcal{A}}|_{U \sqcup V} = \int_{C_2(U) \sqcup C_2(V)} \mathcal{A}^{\boxtimes 2} \otimes \Omega^1_{\log} \simeq \barcomplex{1}{\mathcal{A}}|_U \otimes \barcomplex{1}{\mathcal{A}}|_V\]

\emph{Translation.}
\begin{itemize}
\item $j_{!*}$ (pushforward) = configuration space integral
\item Factorization isomorphism = bar complex splitting
\item Add explicit forms $\Omega^1_{\log}$
\end{itemize}
\end{example}

\subsection{Sign rules for this manuscript}

For easy reference, here are ALL sign rules used in this manuscript in one place:

\subsubsection{Koszul signs}

\emph{Rule.} When permuting two graded objects $a$ and $b$:
\[a \otimes b = (-1)^{|a| \cdot |b|} b \otimes a\]

\emph{Application.} Moving a form $\omega$ of degree $|\omega| = n$ past fields $\phi_1, \ldots, \phi_k$:
\[(\phi_1 \otimes \cdots \otimes \phi_k) \otimes \omega = (-1)^{n(|\phi_1| + \cdots + |\phi_k|)} \omega \otimes (\phi_1 \otimes \cdots \otimes \phi_k)\]

\subsubsection{Collision divisor signs}

\emph{Rule.} Collision divisors are always written with indices in increasing order: $\colldiv{i}{j}$ means $i < j$.

\emph{Convention.} We write $\colldiv{i}{j}$ with $i < j$.  Geometrically $D_{ij} = D_{ji}$ (same subvariety $z_i = z_j$).  The sign convention $\colldiv{j}{i} = -\colldiv{i}{j}$ applies to the \emph{oriented divisor} (i.e., to the associated residue forms, not the subvariety itself).

\subsubsection{Arnold relation signs}

\emph{3-term Arnold relation.}
\[\LogForm{i}{j} \wedge \LogForm{j}{k} + \LogForm{j}{k} \wedge \LogForm{k}{i} + \LogForm{k}{i} \wedge \LogForm{i}{j} = 0\]

\emph{Higher-degree consequences.}
All higher-degree relations in the OS algebra follow from the 3-term Arnold relations by repeated application. There is no independent ``general Arnold relation'' beyond the 3-term version.

\subsubsection{Residue signs}

The sign arising from the residue at $\colldiv{i}{j}$ is determined by the Koszul sign rule: the residue operation (cohomological degree~$1$) must commute past the preceding tensor factors.  For elements $s^{-1}a_1 \otimes \cdots \otimes s^{-1}a_n$ in the bar complex, the residue at $D_{ij}$ (colliding factors $i$ and $j$) picks up the sign $(-1)^{\sum_{k<i}|s^{-1}a_k|}$.  In particular, when all generators have even desuspended degree, all residue signs are positive.

For the explicit sign computations relevant to the bar differential, see the proof of $d^2 = 0$ in \S\ref{sec:bar-nilpotency}.

\subsection{Common pitfalls and how to avoid them}

\subsubsection{Pitfall 1: forgetting Koszul signs}

\emph{Wrong.} $d(a \otimes b \otimes \omega) = da \otimes b \otimes \omega + a \otimes db \otimes \omega$

\emph{Right.} $d(a \otimes b \otimes \omega) = da \otimes b \otimes \omega + (-1)^{|a|} a \otimes db \otimes \omega$

The sign $(-1)^{|a|}$ comes from moving $d$ (degree +1) past $a$ (degree $|a|$).

\subsubsection{Pitfall 2: confusing hat notations}

\emph{Wrong.} $\widehat{ij}$ could mean ``omit factors $i$ and $j$,'' or ``the set $\{i, j\}$,'' or something else entirely.

\emph{Right.} Use our convention: $\omitpair{i}{j}$ always means ``omit both $i$ and $j$'' with no ambiguity.

\subsubsection{Pitfall 3: collision divisor ordering}

\emph{Convention.} We write collision divisors with $i < j$: $\colldiv{1}{3}$ rather than $\colldiv{3}{1}$.  Since $\colldiv{i}{j} = \{z_i = z_j\}$ is an unordered locus, $\colldiv{3}{1} = \colldiv{1}{3}$ (no sign).  Likewise, the logarithmic forms satisfy $\LogForm{j}{i} = \LogForm{i}{j}$ since $d\log(z_j - z_i) = d\log(-(z_i - z_j)) = d\log(z_i - z_j)$.

The sign convention $e_{ji} = -e_{ij}$ applies only to Orlik--Solomon algebra generators (an algebraic convention for exterior algebra presentations), not to the geometric forms or divisors.

\subsubsection{Pitfall 4: Arnold relation orientation}

\emph{Wrong.} Assuming Arnold relations hold without signs.

\emph{Right.} The cyclic sum includes signs from wedge product antisymmetry:
\[\LogForm{1}{2} \wedge \LogForm{2}{3} + \LogForm{2}{3} \wedge \LogForm{3}{1} + \LogForm{3}{1} \wedge \LogForm{1}{2} = 0\]

Note: Each term has a $+$ sign in the cyclic sum, but the individual forms anticommute.

\subsection{Summary and recommendations}

\emph{For readers unfamiliar with sign conventions.}
\begin{enumerate}
\item Start with Remark~\ref{rem:LV-signs} (Loday--Vallette comparison, in Chapter~\ref{chap:bar-cobar})
\item Consult this appendix when reading other sources
\item Work through the examples in \S\ref{sec:heisenberg-bar-complex} (Heisenberg explicit calculations)
\item Verify signs match in low-degree computations
\end{enumerate}

\emph{For experts.}
Our conventions match:
\begin{itemize}
\item Kontsevich's graph complex (lexicographic ordering)
\item Costello--Gwilliam's factorization algebras (cohomological grading)
\item Beilinson--Drinfeld's D-module perspective (implicit signs)
\end{itemize}

The main difference from Loday--Vallette is grading direction ($|d| = +1$ vs $-1$); the translation is a global sign reversal on the grading.

\subsection{Sign convention dictionary}
\label{sec:sign-dictionary-complete}

\begin{table}[ht!]
\centering
\caption{Sign conventions across sources}
\label{tab:sign-conventions-complete}
\small
\begin{tabular}{|l|c|c|c|c|}
\hline
\textbf{Item} & \textbf{Ours} & \textbf{BD} & \textbf{CG} & \textbf{LV} \\
\hline
\hline
Differential degree & $|d| = +1$ & implicit & $+1$ & $-1$ \\
Koszul sign rule & $(-1)^{|a||b|}$ & yes & yes & yes \\
Suspension & $s: V \to V[-1]$ & yes & yes & opposite \\
Bar complex degree & increasing & yes & yes & decreasing \\
Logarithmic 1-forms & $\eta_{ij} = \frac{dz_i - dz_j}{z_i - z_j}$ & yes & yes & N/A \\
Residue sign & $\text{Res}_{z=w}[\eta_{ij}] = +1$ & yes & yes & N/A \\
Factorization & left-to-right & yes & opposite & N/A \\
Collision divisor & $D_{ij}$ with $i < j$ & yes & same & N/A \\
Transposition & $\eta_{ji} = \eta_{ij}$ (forms) & yes & yes & N/A \\
Wedge product & $a \wedge b = (-1)^{|a||b|} b \wedge a$ & yes & yes & yes \\
Verdier duality & $\mathbb{D}(\mathcal{F}) = R\mathcal{H}om(\mathcal{F}, \omega[d])$ & yes & same & N/A \\
Central extension & $[a,b] = ab - ba$ & yes & yes & yes \\
A-infinity & $\mu_n: \mathcal{A}^{\otimes n} \to \mathcal{A}$ & implicit & yes & yes \\
\hline
\end{tabular}
\end{table}

\emph{Key to abbreviations.}
\begin{itemize}
\item BD = Beilinson--Drinfeld \cite{BD04}
\item CG = Costello--Gwilliam \cite{CG17}
\item LV = Loday--Vallette \cite{LV12}
\end{itemize}

\subsubsection{Conversion formulas}

\begin{proposition}[Loday--Vallette conversion; \ClaimStatusProvedHere]\label{prop:LV-conversion-complete}
To convert from our conventions to Loday--Vallette:
\begin{enumerate}
\item Replace $|d| = +1$ with $|d| = -1$ (flip grading direction)
\item Replace bar complex $\bar{B}^n$ with $\bar{B}_{-n}$ (decreasing filtration)
\item Keep Koszul sign rule unchanged (all four conventions agree on the Koszul sign rule)
\item Keep composition order unchanged: $(f \circ g)_{\text{LV}} = f \circ g$ (the degree reversal affects signs in the Leibniz rule, not the order of function composition)
\end{enumerate}
\end{proposition}

\begin{proof}[Verification]
The key identity is that homological grading (LV) relates to cohomological grading
(ours) by $\deg_{\text{hom}} = -\deg_{\text{coh}}$.

For a differential $d$:
\begin{itemize}
\item Our convention: $d: C^n \to C^{n+1}$, so $|d| = +1$
\item LV convention: $d: C_n \to C_{n-1}$, so $|d|_{\text{LV}} = -1$
\end{itemize}

The Koszul sign rule in the Leibniz formula is the same in both conventions:
\[d(a \otimes b) = da \otimes b + (-1)^{|a|} a \otimes db\]
(The degree $|a|$ is taken in whichever grading is active.)  The sign difference between conventions manifests in the \emph{bar differential}, where LV uses explicitly suspended elements $sa$: the bar differential sign $(-1)^{|sa|} = (-1)^{|a|+1}$ in LV's homological convention (where $|sa| = |a|+1$) differs from $(-1)^{|a|}$ in ours.
\end{proof}

\subsubsection{Explicit sign calculations}

\begin{example}[Three-point Arnold relation]\label{ex:arnold-sign-explicit-dict}
For three points $z_1, z_2, z_3$, the Arnold relation is:
\[\eta_{12} \wedge \eta_{23} + \eta_{23} \wedge \eta_{31} + \eta_{31} \wedge \eta_{12} = 0\]

\emph{Sign verification.}
\begin{enumerate}
\item Each $\eta_{ij}$ has degree 1
\item Wedge product: $\eta \wedge \eta'$ has degree 2
\item Moving forms: $\eta \wedge \eta' = (-1)^{1 \cdot 1} \eta' \wedge \eta = -\eta' \wedge \eta$
\end{enumerate}

All signs are explicit and match BD conventions.
\end{example}

\begin{remark}[Fermionic vs bosonic]\label{rem:fermionic-bosonic-signs-dict}
For fermionic fields $\psi$: $\psi(z)\psi(w) = -\psi(w)\psi(z)$ (Koszul sign with $|\psi| = 1$).

For bosonic fields $\phi$: $\phi(z)\phi(w) = \phi(w)\phi(z)$ (degree $|\phi| = 0$).

The bar complex treats both uniformly via suspension $s: V \to V[-1]$.
\end{remark}


\subsection{Master comparison table}
\label{sec:master-comparison-signs}

\begin{table}[H]
\centering
\caption{Sign conventions across four sources}
\renewcommand{\arraystretch}{1.6}
\footnotesize
\begin{tabular}{|l|c|c|c|c|}
\hline
\textbf{Convention} & \textbf{BD} & \textbf{CG} & \textbf{LV} & \textbf{Us} \\
\hline
\hline
\textbf{Koszul sign} & $(-1)^{|a||b|}$ & $(-1)^{|a||b|}$ & $(-1)^{|a||b|}$ & $(-1)^{|a||b|}$ \\
\hline
\textbf{Bar differential} & $d + d'$ & $d_{\text{int}} + d_{\text{ext}}$ & $d_1$ & $d_{\text{int}} + d_{\text{res}}$ \\
\hline
\textbf{Cobar degree} & $[n]$ & $[n-2]$ & $[-n]$ & $[2-n]$ \\
\hline
\textbf{Suspension} & $s: V \to V[-1]$ & $s: V \to V[-1]$ & $s: V \to V[1]$ & $s: V \to V[-1]$ \\
\hline
\textbf{Desuspension} & $s^{-1}: V \to V[1]$ & $s^{-1}: V \to V[1]$ & $s^{-1}: V \to V[-1]$ & $s^{-1}: V \to V[1]$ \\
\hline
\textbf{Twist map} & $(-1)^{|a||b|}(b \otimes a)$ & $(-1)^{|a||b|}(b \otimes a)$ & $(-1)^{|a||b|}(b \otimes a)$ & $(-1)^{|a||b|}(b \otimes a)$ \\
\hline
\textbf{Differential degree} & $+1$ & $+1$ & $-1$ & $+1$ \\
\hline
\textbf{Coalgebra comult} & $\Delta$ & $\Delta$ & $\Delta^!$ & $\Delta$ \\
\hline
\textbf{Residue orientation} & counterclockwise & contour dep. & N/A & counterclockwise \\
\hline
\textbf{$\hbar$ convention} & not used & $\hbar = 1$ & not used & $\hbar$ explicit \\
\hline
\end{tabular}
\end{table}

\subsubsection{Detailed conversion formulas}

\paragraph{Koszul signs.}

\begin{definition}[Koszul sign rule]
When permuting graded elements $a$ and $b$ in a tensor product:

\emph{BD, CG, Us convention.}
\[a \otimes b \leftrightarrow b \otimes a \quad \text{with sign} \quad (-1)^{|a| \cdot |b|}\]

\emph{LV convention.}
\[a \otimes b \leftrightarrow b \otimes a \quad \text{with sign} \quad (-1)^{|a| \cdot |b|} \quad \text{(same Koszul rule)}\]

\emph{Key difference.}
LV works in homological grading ($|d| = -1$), using suspended elements $sa$ with $|sa| = |a| + 1$. The apparent sign differences in bar complex formulas arise from the shifted degrees of $sa$, not from a different Koszul rule. Specifically, swapping $sa \otimes sb$ gives $(-1)^{|sa||sb|} = (-1)^{(|a|+1)(|b|+1)}$, which differs from $(-1)^{|a||b|}$ by the factor $(-1)^{|a|+|b|+1}$.
\end{definition}

\begin{example}[Koszul sign in bar construction]
For $a \in \mathcal{A}^m$, $b \in \mathcal{A}^n$, $c \in \mathcal{A}^p$:

\emph{Us/BD/CG.}
\[d(a \otimes b \otimes c) = d(a) \otimes b \otimes c + (-1)^m a \otimes d(b) \otimes c
+ (-1)^{m+n} a \otimes b \otimes d(c)\]

\emph{LV (on suspended elements $sa, sb, sc$).}
\[d(sa \otimes sb \otimes sc) = d(sa) \otimes sb \otimes sc + (-1)^{m+1} sa \otimes d(sb) \otimes sc
+ (-1)^{m+n+2} sa \otimes sb \otimes d(sc)\]

\emph{Explanation.} The sign differences arise because $|sa| = m+1$, $|sb| = n+1$, $|sc| = p+1$, and the standard Koszul rule is applied to these shifted degrees. The LV Koszul sign rule itself is the same $(-1)^{|x||y|}$.
\end{example}

\subsubsection{Quick translation table}

\begin{table}[H]
\centering
\caption{Translation formulas between conventions}
\begin{tabular}{|l|c|}
\hline
\textbf{To convert from} & \textbf{Apply this transformation} \\
\hline
\hline
LV $\to$ Us/BD/CG & Multiply internal terms by $(-1)^{|a|}$, add chiral product \\
\hline
CG $\to$ Us & Shift notation: CG's $V[1]$ (convention $V[n]^k = V^{k-n}$) equals our $V[-1]$ \\
\hline
BD $\to$ Us & Split $d''$ into $d_{\text{residue}}$ and $d_{\text{correction}}$ \\
\hline
Us $\to$ BD & Combine $d_{\text{residue}} + d_{\text{correction}}$ into $d''$ \\
\hline
Us $\to$ LV & Remove chiral product, multiply by $(-1)^{|a|}$ \\
\hline
\end{tabular}
\end{table}

\subsubsection{Recommendations for readers}

\begin{itemize}
\item \emph{Reading BD.} Our conventions match closely; main difference is we separate out
quantum corrections explicitly.
\item \emph{Reading CG.} The conventions match exactly; use their formulas directly.
\item \emph{Reading LV.} Add extra $(-1)^{|a|}$ signs and remember to include chiral product
which LV does not have.
\item \emph{Writing papers.} We recommend the CG conventions (used here) as they are most explicit about quantum corrections.
\end{itemize}

\begin{remark}[Origins of different conventions]
The conventions reflect different origins: BD developed for D-modules (orientation from holomorphic structure), CG optimized for BV formalism and factorization algebras, LV for pure operadic theory (no geometric residues), and ours emphasizing geometric realization and explicit calculations.
All conventions are \emph{mathematically equivalent}---they are different presentations of the same underlying structures.
\end{remark}
